% F085｜八类模型的简短 LaTeX 结论：按四个代数 + 四个三角组织。
\textbf{F085｜常见抽象函数方程的八类模型识别}

\textbf{解决的问题：}遇到只给函数关系、不告诉解析式的题目，先看“输入怎样运算、函数值怎样运算”，直接用特殊值求解；不必先猜具体函数。

\textbf{第一组：四个代数模型}
\[
\begin{array}{rll}
\text{1. 指数型}&f(x+y)=f(x)f(y)&(a^x),\\[3pt]
\text{2. 对数型}&f(xy)=f(x)+f(y)&(\log_a x),\\[3pt]
\text{3. 幂函数型}&f(xy)=f(x)f(y)&(x^a),\\[3pt]
\text{4. 正比例型}&f(x+y)=f(x)+f(y)&(kx).
\end{array}
\]

\textbf{第二组：四个三角函数模型}
\[
\begin{aligned}
\text{5. 余弦型： }&f(x+y)+f(x-y)=2f(x)f(y),\\[3pt]
\text{6. 正切型： }&f(x+y)=\frac{f(x)+f(y)}{1-f(x)f(y)},\\[5pt]
\text{7. 正弦型： }&f(x+y)f(x-y)=f^2(x)-f^2(y),\\[3pt]
\text{8. 余切型： }&f(x+y)=\frac{f(x)f(y)-1}{f(x)+f(y)}.
\end{aligned}
\]
典型函数依次是 $\cos x$、$\tan x$、$\sin x$、$\cot x$。

\textbf{快记：}加变乘想指数，乘变加想对数，乘变乘想幂，加变加想正比例；和差加想余弦，和差乘想正弦；正切和在上，余切积减一在上。

\textbf{使用条件：}对数、幂型通常在正数上讨论；正切/余切需检查函数定义域和分母非零。\textbf{这些都是识别线索，不是无条件的唯一解析式。}
